% Generated mechanically from results/v34/adaptive_benchmark_v34.csv. Do not edit by hand.
\providecommand{\BenchmarkRotatingError}{\ensuremath{6.019\times 10^{-5}}}
\providecommand{\BenchmarkRotatingStages}{60}
\providecommand{\BenchmarkNearestGuardError}{\ensuremath{6.092\times 10^{-5}}}
\providecommand{\BenchmarkNearestGuardStages}{84}
\providecommand{\BenchmarkFirstGuardBelowError}{\ensuremath{1.902\times 10^{-5}}}
\providecommand{\BenchmarkFirstGuardBelowStages}{84}
\providecommand{\BenchmarkGuardAccepted}{15}
\providecommand{\BenchmarkGuardDirectAccepted}{0}
\providecommand{\BenchmarkGuardProjectedAccepted}{0}
\providecommand{\BenchmarkGuardRotatingAccepted}{15}
\providecommand{\BenchmarkGuardLocatorCalls}{15}
\providecommand{\BenchmarkGuardSuccessfulProjections}{0}
\providecommand{\BenchmarkGuardEmptyComponents}{15}
\providecommand{\BenchmarkGuardUncertainOutcomes}{0}
\providecommand{\BenchmarkGuardRootIsolationOps}{45}
\providecommand{\BenchmarkGuardCertificationOps}{117}
\providecommand{\BenchmarkGuardBinaryRecertifications}{0}
\providecommand{\BenchmarkGuardPrecisionEscalations}{0}
\providecommand{\BenchmarkGuardFallbackEvents}{15}
\providecommand{\BenchmarkBufferedError}{\ensuremath{2.853\times 10^{-4}}}
\providecommand{\BenchmarkBufferedDirect}{3}
\providecommand{\BenchmarkBufferedProjected}{2}
\providecommand{\BenchmarkBufferedLocatorCalls}{2}
\providecommand{\BenchmarkBufferedRootIsolationOps}{6}
\providecommand{\BenchmarkBufferedCertificationOps}{25}
\providecommand{\BenchmarkBufferedBinaryRecertifications}{2}
\providecommand{\BenchmarkBufferedStages}{84}
\providecommand{\BenchmarkNearSaddleError}{\ensuremath{2.103\times 10^{-4}}}
\providecommand{\BenchmarkNearSaddleProjected}{1}
\providecommand{\BenchmarkNearSaddleLocatorCalls}{1}

\begin{frontmatter}

\title{Complete-positivity topology of Runge--Kutta maps for qubit Lindblad equations}

\author[odu]{G. Blake Pierpoint\corref{cor1}}
\ead{pierpogb@odu.edu}
\author[uga,weyl]{Olivier Bernard}
\ead{oli.bernard@weylcenter.org}
\author[sjtu,weyl]{Yichen Liu}
\ead{Yichen.Liu@weylcenter.org}
\cortext[cor1]{Corresponding author.}

\affiliation[odu]{organization={Old Dominion University},addressline={5115 Hampton Boulevard},city={Norfolk},state={Virginia},postcode={23529},country={USA}}
\affiliation[uga]{organization={UFR PhITEM (Physics, Engineering, Earth, Environment, Mechanics), Universit\'e Grenoble Alpes},addressline={230 rue de la Physique},city={Saint-Martin-d'H\`eres},postcode={38400},country={France}}
\affiliation[sjtu]{organization={Department of Computer Science and Engineering, Shanghai Jiao Tong University},addressline={800 Dongchuan Road, Minhang District},city={Shanghai},postcode={200240},country={China}}
\affiliation[weyl]{organization={Weyl Center for Mathematical Physics},addressline={1235 Pennsylvania Avenue SE, Unit 1294},city={Washington},state={DC},postcode={20003},country={USA}}

\begin{abstract}
An exact Lindblad propagator is completely positive and trace preserving (CPTP); a direct Runge--Kutta (RK) step applied to the same generator need not be. For the autonomous phase-covariant qubit generator (coherent precession, bidirectional relaxation, pure dephasing) we derive necessary and sufficient conditions for a direct RK step to be CPTP, expressed through the stability function evaluated at the generator's eigenvalues. For one-way relaxation without dephasing these conditions classify classical RK4 into six admissibility regimes as the precession-to-relaxation ratio varies, including regimes where the admissible step set is disconnected or reduces to isolated points. Exactly removing the commuting precession in a rotating frame eliminates the high-frequency disconnected component, so finite-step physicality depends on the numerical representation, not only on the method. Pointwise certificates at four transverse-field strengths show the disconnected structure persists under a noncommuting perturbation. Candidate construction matters independently: at precession ratio 2 the positive-step CPTP sets of one full RK4 step and of two half steps are disjoint, and the fifth-order Richardson extrapolate of RK4 is non-CPTP on the entire step range where both constituents are CPTP. We implement a tri-state guard that certifies the executed binary64 step and candidate before acceptance and fails closed on unresolved decisions. In adaptive tests, error-only step doubling accepts non-CPTP candidates; the guarded workflows project to a certified step or fall back to an exact subflow. The executed candidate, not the method label or the error estimate, determines finite-step physicality.
\end{abstract}

\begin{keyword}
Lindblad equation \sep complete positivity \sep Runge--Kutta methods \sep candidate maps \sep adaptive time stepping \sep certified numerics
\end{keyword}

\end{frontmatter}

\ifdefined\ReviewMode\linenumbers\fi

\section{Introduction}
\label{sec:introduction}

A finite-dimensional Markovian open quantum system evolves under a Gorini--Kossakowski--Sudarshan--Lindblad (GKSL) generator. Its exact propagator is completely positive and trace preserving (CPTP) at every time \cite{Lindblad1976,GoriniKossakowskiSudarshan1976,BreuerPetruccione2002}. Complete positivity is stronger than trace preservation or positivity on one selected trajectory: the map must remain positive after extension by an arbitrary ancilla, equivalently its Choi matrix must be positive semidefinite \cite{Choi1975CP,Kraus1971,Stinespring1955}. This distinction is operational whenever a numerical update is reused for different inputs or applied to one subsystem of a larger register.

Standard direct RK schemes are not CPTP-preserving in general, and explicit RK formulas supply no generator-independent complete-positivity guarantee \cite{RieschJirauschek2019,AppeloCheng2025}. Numerical approaches to Lindblad evolution have a longer history: high-order stochastic unravelings, deterministic and adaptive low-rank approximations, and locally purified tensor networks predate the current wave of CPTP-by-construction methods \cite{SteinbachGarrawayKnight1995,LeBrisRouchon2013,LeBrisRouchonRoussel2015,WernerEtAl2016}. Recent Kraus-form, quantum-channel, exponential, low-rank, nested-Picard, and tensor-compressed schemes instead alter the update so that positivity or complete positivity is built into the method \cite{RieschPiklJirauschek2020,CaoLu2025,AppeloCheng2025,RobinRouchonSellem2025,ChenEtAl2026,GodinezRamirez2025,HuAppeloCheng2025NPI,DelMastroAppeloCheng2026TT,HuAppeloCheng2026Gregory}. Complementary work addresses deterministic propagation, a posteriori error control, randomized simulation, and optimized unraveling \cite{GuZhang2025,TureHyeonJang2026,EtienneyRobinRouchon2026,ChenLiLuYing2025,CaoHeLi2026Unraveling}.

We ask the complementary question without changing the RK formula: for a fixed generator and a specified numerical candidate construction, which finite-step maps are physical? The answer is not determined by convergence order, scalar absolute stability, componentwise positivity, absolute monotonicity, or strong-stability preservation \cite{BolleyCrouzeix1978,Kraaijevanger1986Absolute,Kraaijevanger1991Contractivity,Horvath1998Positivity,Horvath2005Threshold,GottliebShuTadmor2001,FerracinaSpijker2004,Higueras2005,BlanesIserlesMacnamara2022}. Complete positivity couples distinct Liouvillian modes through a channel inequality, so every scalar mode can be stable while the induced map lies outside the CPTP set.

The exact laboratory is the autonomous phase-covariant qubit family: coherent $z$-precession, downward and upward relaxation, and pure dephasing. The underlying dynamics and non-unital qubit-channel geometry are established \cite{FujiwaraAlgoet1999,RuskaiSzarekWerner2002,BraunEtAl2014,Hall2008Complete,Lankinen2016Complete,Filippov2020Phase,Siudzinska2022Mixtures,Siudzinska2023Geometry,LonigroChruscinski2023}. Pulling that geometry back through an RK stability function produces a one-dimensional semialgebraic CPTP stability set along each generator ray. To our knowledge, prior work has not classified its connected, reentrant, isolated, and detached components for an unchanged RK formula, nor the consequences of constructing several different candidate maps over the same interval.

On the laboratory-frame one-way, zero-dephasing ray, classical RK4 provides a closed-form topology. In the laboratory-frame one-way problem its CPTP set changes repeatedly as $\varpi=|\omega|/\Gamma$ varies. A narrow reentrant band occurs for $0.82736\ldots<\varpi<\sqrt3/2$; above a second bifurcation the positive admissible interval is detached from the origin. At high frequency its dimensional edges satisfy
\begin{equation}
 h_-(\varpi)=\frac{3\Gamma}{\omega^2}+O\!\left(\frac{\Gamma^3}{\omega^4}\right),
 \qquad
 h_+(\varpi)=\frac{2\sqrt2}{|\omega|}+O\!\left(\frac{\Gamma}{\omega^2}\right).
 \label{eq:intro_window}
\end{equation}
This topology is a property of the discretization--representation pair. An exact rotating frame removes the coherent frequency from the commuting benchmark. In the displayed transverse-field family, separate exact calculations at $\Omega_x/\Gamma\in\{0,0.05,0.10,0.20\}$ each retain two map-level CP components; these four pointwise certificates do not establish uniform persistence on the intervening interval.

Adaptive workflows add a second, more general source of structure. Step doubling constructs a coarse full-step map and a fine composition of two half steps; local extrapolation may form a negative-weight affine combination \cite{Shampine1985Doubling,NussleinRanochaKetcheson2021,Szyszkowicz1985Richardson,ZlatevEtAl2017Richardson,HavasiKazemi2018Richardson}. These are distinct linear maps over one interval. At $\varpi=2$ the nontrivial positive-step CPTP components of the full and two-half-step candidates are disjoint, so no positive total step makes both candidates physical. The Richardson extrapolate built from RK4 is an order-five, eighth-degree stability function whose first exponential defect points out of the CPTP set. A global Sturm certificate then proves non-CPTP behavior throughout a positive interval on which both constituent candidates are CPTP. Unlike the frequency-induced laboratory-frame topology, this nonrotating extrapolation failure is unaffected by an exact rotating-frame change.

Direct RK remains relevant because it is widely available, works naturally with sparse or operator-form Liouvillian actions, and is embedded in general scientific ODE workflows \cite{PuzzuoliEtAl2023}. Time-dependent or noncommuting driven generators may not admit one static frame that removes all fast terms at every stage. Frame transformations, integrating factors, splitting, and CPTP-by-construction propagators provide alternative numerical representations. Exact knowledge of the unchanged method's topology informs diagnostics, frame selection, component projection, and fallback decisions.

The exact Choi-block criterion holds for the full autonomous phase-covariant family. The six-regime classification is restricted to the laboratory-frame one-way, zero-dephasing ray. We next test whether the disconnected topology persists under local buffers and a noncommuting transverse field. We then analyze the full-step, equal-substep, mixture, and Richardson candidates, obtaining the disjoint full-step/two-half-step result and the global non-CPTP interval for the RK4 Richardson candidate.

The controller uses three outcomes---\textsc{pass}, \textsc{fail}, and \textsc{uncertain}---and certifies the exact binary64 step together with the construction rule for the candidate actually executed. A rejection caused only by complete positivity leaves the proportional--integral (PI) controller's error memory unchanged. The controller projects in the buffered and near-saddle tests and falls back in the exact-boundary test.

Sections~\ref{sec:criterion}--\ref{sec:rk4topology} derive the criterion and boundary topology. The remaining sections treat robustness, candidate maps, certified step control, operational consequences, and structure-preserving alternatives.

\section{Direct RK propagation and the exact channel criterion}
\label{sec:criterion}

We use units with $\hbar=1$. Let $|0\rangle$ and $|1\rangle$ denote the lower- and upper-energy states, $\sigma_-=|0\rangle\langle1|$, $\sigma_+=|1\rangle\langle0|$, and $\sigma_z=|0\rangle\langle0|-|1\rangle\langle1|$. The autonomous phase-covariant generator is
\begin{equation}
 \mathcal L(\rho)=-\ii[H,\rho]+\gamma_\downarrow\mathcal D[\sigma_-](\rho)
 +\gamma_\uparrow\mathcal D[\sigma_+](\rho)
 +\frac{\gamma_\phi}{2}(\sigma_z\rho\sigma_z-\rho),
 \label{eq:gksl}
\end{equation}
where $H=-\omega\sigma_z/2$, $\mathcal D[L](\rho)=L\rho L^\dagger-\{L^\dagger L,\rho\}/2$, and all rates are nonnegative \cite{Hall2008Complete,Lankinen2016Complete,Filippov2020Phase,Siudzinska2023Geometry}. Define
\begin{equation}
 \Gamma=\gamma_\downarrow+\gamma_\uparrow,
 \quad \theta=\gamma_\uparrow/\Gamma,
 \quad x=\Gamma h,
 \quad u=\gamma_\phi h,
 \quad \nu=\omega h.
 \label{eq:coordinates}
\end{equation}
For $\Gamma>0$, $\theta\in[0,1]$ is the stationary excitation fraction. It represents a thermal stationary population only when the rates satisfy an additional detailed-balance relation. The population and coherence modes obey
\begin{equation}
 \dot q=-\Gamma(q-\theta),
 \qquad
 \dot\rho_{01}=-\left(\frac{\Gamma}{2}+\gamma_\phi-\ii\omega\right)\rho_{01}.
 \label{eq:modes}
\end{equation}
The boundary $\Gamma=0$ is recovered continuously and does not require $\theta$.

For an $s$-stage RK method with real Butcher coefficients $(A,b)$, let
\begin{equation}
 R(z)=1+z b^T(I-zA)^{-1}{\bf 1}
 \label{eq:stability_function}
\end{equation}
be its stability function \cite{HairerNorsettWanner1993}. We retain the true stage-evaluation domain: a removable cancellation in a reduced rational expression does not validate a singular stage system. Direct application to Eq.~\eqref{eq:modes} gives
\begin{equation}
 a=R(-x),
 \qquad
 c=R(-x/2-u+\ii\nu).
 \label{eq:multipliers}
\end{equation}
The resulting complex-linear map on $X=(X_{jk})\in M_2(\mathbb C)$ is
\begin{equation}
 \Phi_{R}^{(\theta)}(X)=
 \begin{pmatrix}
 A_\theta X_{00}+C_\theta X_{11} & cX_{01}\\
 c^*X_{10} & B_\theta X_{00}+D_\theta X_{11}
 \end{pmatrix},
 \label{eq:direct_map}
\end{equation}
where
\begin{align}
 A_\theta&=1-\theta(1-a), & B_\theta&=\theta(1-a),\nonumber\\
 C_\theta&=(1-\theta)(1-a), & D_\theta&=\theta+(1-\theta)a.
 \label{eq:ABCD}
\end{align}
The map preserves trace and Hermiticity exactly. No integrating factor, projection, Kraus reconstruction, or post-step renormalization has been introduced.

Using the convention $J(\Phi)=\sum_{jk}|j\rangle\langle k|\otimes\Phi(|j\rangle\langle k|)$, the unnormalized Choi matrix is
\begin{equation}
 J_\theta=
 \begin{pmatrix}
 A_\theta&0&0&c\\
 0&B_\theta&0&0\\
 0&0&C_\theta&0\\
 c^*&0&0&D_\theta
 \end{pmatrix}.
 \label{eq:thermal_choi}
\end{equation}
This is the established non-unital phase-covariant Choi form \cite{Choi1975CP,RuskaiSzarekWerner2002,BraunEtAl2014,Siudzinska2023Geometry}. Its pullback through Eq.~\eqref{eq:multipliers} yields the exact numerical criterion.

\begin{theorem}[Finite-step CPTP criterion]
\label{thm:criterion}
On the RK stage-evaluation domain, the direct step \eqref{eq:direct_map} is CPTP if and only if
\begin{equation}
 a\le1,
 \qquad A_\theta\ge0,
 \qquad D_\theta\ge0,
 \qquad \Fmargin(x,u,\nu)\ge0,
 \label{eq:full_criterion}
\end{equation}
where
\begin{equation}
 \Fmargin(x,u,\nu)=a+\theta(1-\theta)(1-a)^2-|c|^2.
 \label{eq:full_margin}
\end{equation}
\end{theorem}

\begin{proof}
The scalar Choi entries $B_\theta$ and $C_\theta$ are nonnegative exactly when $a\le1$. The remaining Hermitian $2\times2$ block is positive semidefinite exactly when its diagonal entries and determinant are nonnegative. Direct multiplication gives $A_\theta D_\theta=a+\theta(1-\theta)(1-a)^2$.
\end{proof}

At the one-way endpoint $\theta=0$, Eq.~\eqref{eq:full_criterion} reduces to
\begin{equation}
 0\le a\le1,
 \qquad
 M_R(x,u,\nu):=a-|c|^2\ge0.
 \label{eq:one_way_criterion}
\end{equation}
The margin already implies $a\ge0$. The lower eigenvalue of the active Choi block is
\begin{equation}
 \lambda_-=
 \frac{1+a-\sqrt{(1-a)^2+4|c|^2}}{2}
 =\frac{2M_R}{1+a+\sqrt{(1-a)^2+4|c|^2}},
 \label{eq:lambda_margin}
\end{equation}
so its sign is the sign of $M_R$. The population ceiling and block margin are independent exit mechanisms.

\subsection{Representation dependence under a rotating frame}
\label{subsec:rotating_frame}

The exact physical evolution is invariant under a change of frame, but a specified direct-RK discretization need not be. Let
\begin{equation}
 \widetilde\rho(t)=\e^{\ii Ht}\rho(t)\e^{-\ii Ht}.
\end{equation}
Because $\e^{\ii Ht}\sigma_\pm\e^{-\ii Ht}=\e^{\mp\ii\omega t}\sigma_\pm$ and $\mathcal D[\e^{\ii\phi}L]=\mathcal D[L]$, the transformed generator contains the same dissipators and no coherent $z$-precession. Denote that dissipative generator by $\mathcal L_{\rm dis}$ and the exact back-rotation by the unitary channel $\mathcal U_h$.

\begin{proposition}[Frame-transformed RK step]
\label{prop:rotating_frame}
The numerical procedure
\begin{equation}
 \Phi_h^{\rm RF}=\mathcal U_h\circ R(h\mathcal L_{\rm dis})
 \label{eq:rotating_frame_step}
\end{equation}
is CPTP if and only if $R(h\mathcal L_{\rm dis})$ is CPTP. Its admissibility set is therefore independent of $\omega$. For one-way damping without dephasing, classical RK4 has the attached interval $0\le x\le\alpha_4$ in this frame, rather than the high-frequency detached regions found for direct laboratory-frame RK4.
\end{proposition}

Thus the regions below characterize direct RK applied in the stated representation. This is analogous to the familiar representation dependence of stiffness and phase error: a rotating frame can remove the high-frequency pathology in this benchmark, whereas the noncommuting perturbation studied in Sec.~\ref{sec:robustness} cannot be removed by the same transformation.

\section{Interior buffers and boundary orientation}
\label{sec:local}

The exact semigroup has $a_{\rm ex}=\e^{-x}$ and $|c_{\rm ex}|^2=\e^{-x-2u}$, hence
\begin{equation}
 F_{{\rm ex},\theta}=\e^{-x}(1-\e^{-2u})
 +\theta(1-\theta)(1-\e^{-x})^2.
 \label{eq:exact_buffer}
\end{equation}
Positive dephasing contributes at first order, while bidirectional relaxation contributes at second order. The one-way, zero-dephasing endpoints are singular because both buffers vanish.

\subsection{Local interior laws}

Let $R(z)=1+z+O(z^2)$ be consistent. As $(x,u,\nu)\to0$,
\begin{equation}
 \Fmargin(x,u,\nu)
 =2u+O(x^2+xu+u^2+\nu^2).
 \label{eq:dephasing_buffer}
\end{equation}
Along a fixed generator ray
\begin{equation}
 u=\kappa x,
 \qquad \nu=\varpi x,
 \qquad \kappa=\gamma_\phi/\Gamma,
 \qquad \varpi=\omega/\Gamma,
 \label{eq:generator_ray}
\end{equation}
any $\kappa>0$ supplies a nonzero local CPTP interval.

When $u=0$, write $R(z)=1+z+r_2z^2+O(z^3)$. Expansion of Eq.~\eqref{eq:full_margin} gives
\begin{equation}
 \Fmargin(x,0,\varpi x)
 =C_{R,\theta}(\varpi)x^2+O(x^3),
 \label{eq:excitation_expansion}
\end{equation}
where
\begin{equation}
 C_{R,\theta}(\varpi)
 =\theta(1-\theta)+\frac{r_2}{2}-\frac14+(2r_2-1)\varpi^2.
 \label{eq:excitation_coeff}
\end{equation}
Every method of order at least two has $r_2=1/2$, so
\begin{equation}
 \Fmargin(x,0,\varpi x)
 =\theta(1-\theta)x^2+O(x^3)>0
 \label{eq:excitation_buffer}
\end{equation}
for $0<\theta<1$ and sufficiently small $x$. Dephasing and bidirectional relaxation move the exact trajectory into the relative interior of the relevant channel face. The signed boundary theory below therefore applies precisely at $\theta\in\{0,1\}$ and $u=0$.

On the pure-dephasing/precession boundary $x=0$, the criterion is
\begin{equation}
 |R(-u+\ii\nu)|\le1.
 \label{eq:pure_dephasing}
\end{equation}
The purely unitary slice is governed by imaginary-axis absolute stability, $|R(\ii\nu)|\le1$ \cite{HairerNorsettWanner1993,IserlesNorsett1991OrderStars}.

For later controller estimates, specialize to RK4 and a rotation-dominated joint limit in which $\kappa\to0$ or $\theta(1-\theta)\to0$, $|\varpi|\to\infty$, and the resulting attached endpoint satisfies $x\to0$ and $|\varpi|x\ll1$. The outward boundary defect is
\begin{equation}
 M_4(x;\varpi)=-\frac{\varpi^4x^5}{24}+o(\varpi^4x^5).
\end{equation}
Balancing it against the local buffers gives
\begin{equation}
 x_{\rm att}^{(\phi)}\sim\frac{(48\kappa)^{1/4}}{|\varpi|},
 \qquad
 x_{\rm att}^{(\theta)}\sim
 \left[\frac{24\theta(1-\theta)}{\varpi^4}\right]^{1/3}.
 \label{eq:attached_scaling}
\end{equation}
These are joint-limit estimates, not finite-frequency identities.

\subsection{Signed first-defect law}

Fix the one-way, zero-dephasing boundary and set $\alpha=1/2-\ii\varpi$. Suppose
\begin{equation}
 R(z)-\e^z=\eta_m z^m+O(z^{m+1}),
 \qquad m\ge2,
 \qquad \eta_m\ne0.
 \label{eq:first_defect}
\end{equation}
Then $|\e^{-\alpha x}|^2=\e^{-x}$ cancels all lower-order semigroup terms.

\begin{theorem}[Frequency-dependent boundary law]
\label{thm:defect}
Under Eq.~\eqref{eq:first_defect},
\begin{equation}
 \Mmargin(x)
 =(-1)^m\eta_mG_m(\varpi)x^m+O(x^{m+1}),
 \label{eq:defect_law}
\end{equation}
where
\begin{equation}
 G_m(\varpi)=1-2\operatorname{Re}(1/2-\ii\varpi)^m.
 \label{eq:G_definition}
\end{equation}
For sufficiently small $x>0$, the sign of $(-1)^m\eta_mG_m$ determines whether the RK map enters or exits the channel set, except when $G_m=0$.
\end{theorem}

Define
\begin{equation}
 q=1+4\varpi^2,
 \qquad
 \sigma_m=2^m(\alpha^m+\bar\alpha^m).
 \label{eq:sigma_definition}
\end{equation}
Because $\alpha$ and $\bar\alpha$ solve $z^2-z+q/4=0$,
\begin{equation}
 \sigma_0=\sigma_1=2,
 \qquad
 \sigma_m=2\sigma_{m-1}-q\sigma_{m-2},
 \qquad
 G_m=1-\frac{\sigma_m}{2^m}.
 \label{eq:G_recurrence}
\end{equation}
For the methods used below,
\begin{align}
 G_4&=-\frac18q(q-8),
 &G_5&=-\frac5{16}q(q-4),\nonumber\\
 G_6&=\frac1{32}q(q^2-18q+48),
 &G_7&=\frac7{64}q(q-4)^2.
 \label{eq:G_factors}
\end{align}
At $q=4$, or $\varpi=\sqrt3/2$, $\alpha=\e^{-\ii\pi/3}$ and
\begin{equation}
 G_m(\sqrt3/2)=1-2\cos(m\pi/3).
 \label{eq:unit_circle_rule}
\end{equation}
Thus the unit-circle cancellation occurs when $m\equiv\pm1\pmod6$. A simple zero reverses orientation; an even-multiplicity zero produces tangency.

\begin{proposition}[Degenerate frequency]
\label{prop:degenerate}
If $R(z)-\e^z=\eta_mz^m+\eta_{m+1}z^{m+1}+O(z^{m+2})$ and $G_m(\varpi)=0$, then
\begin{align}
 M_{R,\varpi}(x)
 ={}&\Big[(-1)^{m+1}\eta_{m+1}G_{m+1}(\varpi)\nonumber\\
 &+2(-1)^m\eta_m\operatorname{Re}(\bar\alpha\alpha^m)\Big]x^{m+1}
 +O(x^{m+2}).
 \label{eq:degenerate_coefficient}
\end{align}
\end{proposition}

For SSPRK(3,3) and classical RK4 the cancellation cases factor exactly:
\begin{align}
 M_{{\rm SSP3},\sqrt7/2}(x)&=\frac{x^5(3-2x)}{9},\nonumber\\
 M_{{\rm RK4},\sqrt3/2}(x)&=-\frac{x^6(x-2)^2}{576}.
 \label{eq:threshold_factors}
\end{align}
The principal Dormand--Prince fifth-order formula \cite{DormandPrince1980} has two reversal frequencies,
\begin{equation}
 \varpi=\sqrt{2-\frac{\sqrt{33}}4},
 \qquad
 \varpi=\sqrt{2+\frac{\sqrt{33}}4},
 \label{eq:dp5_thresholds}
\end{equation}
so its small-step orientation is inward, outward, then inward. Formal RK order alone cannot predict complete positivity.

\section{Complete RK4 topology on the laboratory-frame one-way ray}
\label{sec:rk4topology}

Set $\theta=u=0$ and $\nu=\varpi x$. For classical RK4,
\begin{equation}
 R_4(z)=1+z+\frac{z^2}{2}+\frac{z^3}{6}+\frac{z^4}{24}.
\end{equation}
The population ceiling $R_4(-x)\le1$ first becomes active at
\begin{equation}
 \alpha_4=2.785293563405\ldots,
 \qquad
 \alpha_4^3-4\alpha_4^2+12\alpha_4-24=0.
 \label{eq:alpha4}
\end{equation}
The Choi margin admits a closed factorization for every $\varpi$.

\begin{proposition}[Laboratory-frame RK4 margin]
\label{prop:rk4factor}
Let $q=1+4\varpi^2$ and $y=qx$. Then
\begin{equation}
 M_4(x;\varpi)=-\frac{x^5q}{147456}\,\mathcal C_q(qx),
 \label{eq:rk4_factor}
\end{equation}
where
\begin{equation}
 \mathcal C_q(y)=y^3-16y^2-32(q-6)y+384(q-4).
 \label{eq:rk4_cubic}
\end{equation}
Its discriminant is
\begin{equation}
 \operatorname{disc}_y\mathcal C_q
 =16384(q-4)(8q^2-123q+348).
 \label{eq:cubic_discriminant}
\end{equation}
\end{proposition}

Define
\begin{equation}
 q_\pm=\frac{123\pm11\sqrt{33}}{16},
 \qquad
 \varpi_\pm=\frac{\sqrt{107\pm11\sqrt{33}}}{8}.
 \label{eq:qcrit}
\end{equation}
Numerically, $\varpi_-=0.82736225145\ldots$ and $\varpi_+=1.63071202389\ldots$. The critical cubics factor as
\begin{align}
 \mathcal C_{q_-}(y)&=(y-9+\sqrt{33})^2(y-2\sqrt{33}+2),\nonumber\\
 \mathcal C_{4}(y)&=y(y-8)^2,\nonumber\\
 \mathcal C_{q_+}(y)&=(y-9-\sqrt{33})^2(y+2+2\sqrt{33}).
 \label{eq:critical_cubics}
\end{align}

\begin{theorem}[Complete laboratory-frame RK4 classification]
\label{thm:rk4topology}
Let $\mathcal A_4(\varpi)$ be the set of $x\ge0$ for which direct laboratory-frame RK4 is CPTP on the one-way, zero-dephasing ray. Then:
\begin{enumerate}
 \item For $1\le q<q_-$, $\mathcal C_q$ has one positive root $y_3$ and
 \[
  \mathcal A_4=[0,\min\{\alpha_4,y_3/q\}].
 \]
 At $q=q_-$ the same formula holds; the additional double root is an internal tangency.
 \item For $q_-<q<4$, $\mathcal C_q$ has three positive roots $y_1<y_2<y_3$ and
 \[
  \mathcal A_4=[0,y_1/q]\cup[y_2/q,y_3/q].
 \]
 \item At $q=4$, $\mathcal A_4=\{0,2\}$.
 \item For $4<q<q_+$, $\mathcal A_4=\{0\}$.
 \item At $q=q_+$,
 \[
  \mathcal A_4=\left\{0,\frac{9+\sqrt{33}}{q_+}\right\}.
 \]
 \item For $q>q_+$, $\mathcal C_q$ has one negative and two positive roots $y_2<y_3$, and
 \[
  \mathcal A_4=\{0\}\cup[y_2/q,y_3/q].
 \]
\end{enumerate}
\end{theorem}

\begin{proof}
For $x>0$, Eq.~\eqref{eq:rk4_factor} makes the block condition equivalent to $\mathcal C_q(qx)\le0$. The sign at the origin is set by $384(q-4)$, and the number and multiplicity of real roots follow from Eq.~\eqref{eq:cubic_discriminant}. Root sums and products determine their signs. Equation~\eqref{eq:critical_cubics} resolves the equality cases. Exact isolation verifies the population ceiling on the attached low-frequency branch and confirms that all detached endpoints lie below $\alpha_4$. Complete sign charts are supplied in the Supplementary Material.
\end{proof}

Figure~\ref{fig:rk4_phase} shows two forms of reentrance. Just below $\sqrt3/2$, a detached interval coexists with the attached small-step interval. The band is only $0.03866\ldots$ wide in $\varpi$, so a coarse frequency grid can miss it. At $\sqrt3/2$ both components collapse to points. Positive-step CPTP then disappears before a remote interval is born at $\varpi_+$.

The equality cases are exact mathematical members of the CPTP set, but isolated zero-margin points are not robust controller targets. The birth at $\varpi_+$ is a saddle node: the two positive roots separate as $O((\varpi-\varpi_+)^{1/2})$. Companion-matrix roots are consequently ill-conditioned there. All classifications use exact Sturm isolation or certified sign continuation; floating companion roots seed plots only.

\begin{figure*}[t]
 \centering
 \makebox[\textwidth][c]{\includegraphics[width=\JCPFigureWidth]{Fig1_RK4_phase_diagram.pdf}}
 \caption{Exact RK4 topology and conditioning on the laboratory-frame one-way, zero-dephasing ray, with $x=\Gamma h$ and $\varpi=|\omega|/\Gamma$. (a) Attached and detached CPTP components constructed from the cubic branches and population ceiling. The inset resolves the narrow reentrant band; isolated equality points are retained. (b) The detached-window width has the square-root scaling expected at the double-root birth, making uncertified floating root decisions unreliable close to $\varpi_+$.}
 \label{fig:rk4_phase}
\end{figure*}

For $q\to\infty$, the lower positive root balances the $q$-dependent constant and linear terms in Eq.~\eqref{eq:rk4_cubic}, while the upper root scales as $\sqrt q$. Hence
\begin{equation}
 x_-(\varpi)=\frac{3}{\varpi^2}+O(\varpi^{-4}),
 \qquad
 x_+(\varpi)=\frac{2\sqrt2}{|\varpi|}+O(\varpi^{-2}).
 \label{eq:window_asymptotics}
\end{equation}
The upper limit is the classical RK4 imaginary-axis boundary because $|R_4(\ii\nu)|^2-1=\nu^6(\nu^2-8)/576$. The lower edge has no scalar analogue; it is produced by the coupled population--coherence Choi condition.

\section{Persistence beyond the critical phase-covariant boundary}
\label{sec:robustness}

The detached intervals of Theorem~\ref{thm:rk4topology} were derived at $\theta=\kappa=0$, where the exact semigroup rides the Choi determinant boundary. Positive dephasing or bidirectional relaxation restores a small-step buffer. The remote component and the intervening nonphysical gap can nevertheless persist.

\begin{theorem}[Relative-open persistence of disconnectedness]
\label{thm:persistence}
Fix an RK method of order at least two, a frequency ratio $\varpi$, and a compact step interval $K$ contained in the RK stage-evaluation domain. Suppose that at $(\theta,\kappa)=(0,0)$:
\begin{enumerate}
 \item the Choi margin has a detached positive component $(r_2,r_3)\subset K$ bounded by simple roots;
 \item there are points $s_-\in(0,r_2)$, $s_+\in(r_2,r_3)$, and $s_{\rm out}>r_3$ in $K$ with margin signs $-$, $+$, and $-$;
 \item every inactive population and Choi-diagonal constraint has uniform strict slack on closed neighborhoods of these points and roots.
\end{enumerate}
Then a relative-open neighborhood of $(0,0)$ in the physical quadrant $\theta,\kappa\ge0$ has at least two positive-step CPTP components whenever $\theta+\kappa>0$.
\end{theorem}

\begin{proof}
Choose disjoint compact neighborhoods of the three sign samples and of the two simple roots, all lying inside $K$. Compactness and the strict-slack hypothesis give one positive lower bound for every inactive constraint on their union. Continuity of the RK map, Choi entries, and active margin in $(x,\theta,\kappa)$ preserves these inequalities and the three strict margin signs in one common relative-open parameter neighborhood. The implicit-function theorem continues the simple roots near $r_2$ and $r_3$. Equations~\eqref{eq:dephasing_buffer} and \eqref{eq:excitation_buffer} supply an attached positive interval near $x=0$ for every nonzero buffer. The preserved $-,+,-$ samples then force a gap and a remote positive component. The theorem establishes persistence of disconnectedness, not persistence of the boundary absorbing interval below $r_2$.
\end{proof}

For RK4 at $\varpi=2$, exact rational data give the representative classifications in Table~\ref{tab:robust_cases}. Every displayed endpoint is isolated by a Sturm count and every intervening sign is evaluated at an exact rational sample.

\begin{table*}[t]
\centering
\caption{Certified RK4 regions at $\varpi=2$ with positive dephasing and/or excitation fraction. Here $\kappa=\gamma_\phi/\Gamma$. All listed roots are simple.}
\label{tab:robust_cases}
\small
\begingroup
\setlength{\tabcolsep}{9pt}
\renewcommand{\arraystretch}{1.10}
\begin{tabular}{@{}ccccc l@{}}
\toprule
$\theta$ & $\kappa$ & $r_1$ & $r_2$ & $r_3$ & Complete CPTP set \\
\midrule
$0$ & $10^{-3}$ & $0.25334887$ & $0.74069600$ & $1.26954271$ & $[0,r_1]\cup[r_2,r_3]$ \\
$10^{-3}$ & $0$ & $0.12226890$ & $0.74255831$ & $1.27054437$ & $[0,r_1]\cup[r_2,r_3]$ \\
$10^{-3}$ & $10^{-3}$ & $0.26203333$ & $0.73836463$ & $1.26975357$ & $[0,r_1]\cup[r_2,r_3]$ \\
$0$ & $10^{-2}$ & $0.51930822$ & $0.68162161$ & $1.26228015$ & $[0,r_1]\cup[r_2,r_3]$ \\
$10^{-2}$ & $0$ & $0.27693028$ & $0.72089681$ & $1.27239139$ & $[0,r_1]\cup[r_2,r_3]$ \\
\bottomrule
\end{tabular}
\endgroup
\end{table*}

The local estimates \eqref{eq:attached_scaling} predict the order of the first attached endpoint. A multiplicative reducer with factor $\rho\in(0,1)$ starting from $x_0>x_{\rm att}$ reaches that component after approximately
\begin{equation}
 j_{\rm att}=\left\lceil\frac{\log(x_{\rm att}/x_0)}{\log\rho}\right\rceil
 \label{eq:attached_reductions}
\end{equation}
reductions. Thus an off-boundary controller may eventually recover under refinement, but the cost can be large when the buffer is weak.

\subsection{Noncommuting Hamiltonian perturbation}

The two-scalar reduction is special to commuting phase-covariant modes, but disconnectedness is not confined to that subfamily. Consider
\begin{equation}
 \mathcal L_{\epsilon}(\rho)=\mathcal L(\rho)-\ii[\Omega_x\sigma_x/2,\rho],
 \label{eq:transverse_generator}
\end{equation}
with the same dissipators as Eq.~\eqref{eq:gksl}. For $\Omega_x\ne0$, phase covariance is broken and the one-step Choi matrix must be formed from the full $4\times4$ RK superoperator.

\begin{proposition}[Persistence in a non-phase-covariant GKSL family]
\label{prop:noncommuting_persistence}
Let $\mathcal L_\eta$ be a continuous finite-dimensional parameterized GKSL family. If four steps $0<h_a<h_g<h_d<h_o$ satisfy strict Choi-eigenvalue signs $+,-,+,-$ for $R_4(h\mathcal L_{\eta_0})$, then the sign pattern, and hence at least two positive-step CP components, persists for all $\eta$ in a sufficiently small neighborhood of $\eta_0$ within that GKSL family.
\end{proposition}

\begin{proof}
The polynomial $R_4(h\mathcal L_\eta)$ and its Choi matrix depend continuously on $(h,\eta)$. The smallest Choi eigenvalue is continuous for Hermitian matrices. Four strict signs therefore remain unchanged in one common neighborhood. The two positive samples are separated by a negative sample, so they cannot belong to one connected CP component.
\end{proof}

At $\Gamma=1$, $\theta=\kappa=10^{-3}$, $\omega/\Gamma=2$, and $\Omega_x/\Gamma=0.1$, exact characteristic-polynomial brackets give the approximate midpoint values
\begin{equation}
 \lambda_{\min}(0.1)\approx9.5156\times10^{-5},\quad
 \lambda_{\min}(0.5)\approx-2.7778\times10^{-3},
\end{equation}
\begin{equation}
 \lambda_{\min}(0.9)\approx5.8959\times10^{-4},\quad
 \lambda_{\min}(1.4)\approx-3.9227\times10^{-1}.
 \label{eq:noncommuting_samples}
\end{equation}
Exact isolation of every principal-minor root yields
\begin{equation}
 \mathcal A_4=[0,0.26357691\ldots]\cup[0.73588381\ldots,1.26847248\ldots].
\end{equation}
We repeat the exact component calculation at $\Omega_x/\Gamma\in\{0,0.05,0.10,0.20\}$. Each of these four parameter values has an attached and a detached CP component. Figure~\ref{fig:robustness} displays only those certified points; no interpolation or uniform interval certificate is asserted. Proposition~\ref{prop:noncommuting_persistence} supplies local persistence around each strict-sign base point, not a generic theorem for arbitrary noncommuting Lindbladians.

\begin{figure*}[t]
 \centering
 \makebox[\textwidth][c]{\includegraphics[width=\JCPFigureWidth]{Fig2_robustness.pdf}}
 \caption{Disconnected RK4 regions at $\varpi=|\omega|/\Gamma=2$, with $x=\Gamma h$. (a) Exact phase-covariant cases with positive dephasing and/or excitation fraction. (b) Discrete component boundaries certified independently at $\Omega_x/\Gamma\in\{0,0.05,0.10,0.20\}$. Markers are not connected or interpolated; the panel does not certify every coupling between the displayed values.}
 \label{fig:robustness}
\end{figure*}

\section{Candidate-map admissibility}
\label{sec:candidates}

A numerical workflow may construct several approximations over one total interval. Complete positivity belongs to the particular linear map that is advanced. We therefore separate candidate construction from error estimation and from the nominal RK formula.

Write $\Phi_\theta(a,c)$ for the phase-covariant map in Eq.~\eqref{eq:direct_map}; the maps below share the same stationary population parameter $\theta$.

\begin{proposition}[Composition calculus]
\label{prop:composition_calculus}
For two phase-covariant maps with the same $\theta$,
\begin{equation}
 \Phi_\theta(a_2,c_2)\circ\Phi_\theta(a_1,c_1)
 =\Phi_\theta(a_2a_1,c_2c_1).
 \label{eq:composition_calculus}
\end{equation}
Consequently, $n$ equal RK substeps over total coordinates $(x,u,\nu)$ have multipliers
\begin{equation}
 a_n=R(-x/n)^n,
 \qquad
 c_n=R[-(x/2+u-\ii\nu)/n]^n.
 \label{eq:n_substep_multipliers}
\end{equation}
\end{proposition}

\begin{proof}
Substitution into Eq.~\eqref{eq:direct_map} shows that the population deviation from the common fixed point is multiplied by $a_2a_1$, while the coherence is multiplied by $c_2c_1$.
\end{proof}

Let $\mathcal A_R$ be the full one-step CPTP set on a fixed one-way generator ray. If the substep population multiplier $a_0=R(-x/n)$ is nonnegative, then $a_0^n\le1$ exactly when $a_0\le1$, and $a_0^n-|c_0|^{2n}$ has the same sign as $a_0-|c_0|^2$. Hence
\begin{equation}
 \Phi_R(x/n)^n\ \text{is CPTP}
 \quad\Longleftrightarrow\quad
 \Phi_R(x/n)\ \text{is CPTP}.
 \label{eq:no_repair_substep}
\end{equation}
Classical RK4 satisfies $R_4(-x)>0$ for every $x\ge0$ because
\begin{equation}
 24R_4(-x)=(x^2-2x)^2+8(x^2-3x+3)>0.
\end{equation}
Hence the total-step admissible set of an $n$-substep RK4 candidate is exactly
\begin{equation}
 \mathcal A_4^{(n)}=n\mathcal A_4.
 \label{eq:dilated_candidate_set}
\end{equation}

\begin{proposition}[Substep-ladder overlap of positive components]
\label{prop:substep_ladder}
Suppose a nontrivial positive-step component is $\mathcal I=[x_-,x_+]$ with $0<x_-<x_+$. Then
\begin{equation}
 \mathcal I^{(n)}\cap\mathcal I^{(n+1)}\ne\varnothing
 \quad\Longleftrightarrow\quad
 \frac{x_+}{x_-}\ge\frac{n+1}{n}.
 \label{eq:ladder_overlap}
\end{equation}
For the full-step/two-half-step pair ($n=1$), the nontrivial positive-step components are disjoint exactly when $x_+/x_-<2$. The same factor-two specialization is equivalent to the existence of initial proposals for which repeated halving skips $\mathcal I$. More generally, a fixed reduction factor $\rho$ hits the component for every proposal above it if and only if $\rho\ge x_-/x_+$.
\end{proposition}

\begin{proof}
Equation~\eqref{eq:dilated_candidate_set} gives the corresponding positive components $\mathcal I^{(n)}=[nx_-,nx_+]$ and $\mathcal I^{(n+1)}=[(n+1)x_-,(n+1)x_+]$. Their overlap is equivalent to $(n+1)x_-\le nx_+$. The reduction statement follows by comparing two consecutive proposals $\rho^jx_0>x_+$ and $\rho^{j+1}x_0<x_-$.
\end{proof}

At $\varpi=2$,
\begin{equation}
 \mathcal A_4=\{0\}\cup[0.7447824278959\ldots,1.2703346269739\ldots],
 \label{eq:direct_varpi2_interval}
\end{equation}
whereas
\begin{equation}
 \mathcal A_4^{(2)}=\{0\}\cup[1.4895648557918\ldots,2.5406692539478\ldots].
 \label{eq:fine_varpi2_interval}
\end{equation}
The sets share only the zero step; their nontrivial positive-step components are disjoint. Thus $H=1$ gives a CPTP full candidate and a non-CPTP fine candidate, whereas $H=2$ reverses those statuses. The step-doubling pair is also the most demanding consecutive pair: at $\varpi=2$, the corresponding positive components $\mathcal I_4^{(n)}$ and $\mathcal I_4^{(n+1)}$ overlap for every $n\ge2$. Positive-component disjointness and the halving-skip condition end together at
\begin{equation}
 \varpi_{1/2}=2.248254604514980\ldots,
 \label{eq:halving_threshold}
\end{equation}
whose exact algebraic certificate is given in the Supplementary Material.

\subsection{Affine candidates and the extrapolation trap}

The CPTP set is convex. If $\Phi_j$ are CPTP, $w_j\ge0$, and $\sum_jw_j=1$, then $\sum_jw_j\Phi_j$ is CPTP. Trace preservation survives every affine combination with weights summing to one, but negative weights need not preserve complete positivity.

For classical RK4, let $\Phi_{\rm c}(x)$ be the coarse full-step map and $\Phi_{\rm f}(x)=\Phi_{\rm RK4}(x/2)^2$ the fine map. The standard local Richardson extrapolate built from fourth-order RK4 is \cite{Szyszkowicz1985Richardson,ZlatevEtAl2017Richardson,HavasiKazemi2018Richardson}
\begin{equation}
 \Phi_{\rm ext}(x)=\frac{16\Phi_{\rm f}(x)-\Phi_{\rm c}(x)}{15}.
 \label{eq:richardson_map}
\end{equation}
The affine candidate is itself generated by the stability polynomial
\begin{equation}
 R_{\rm ext}(z)=\frac{16R_4(z/2)^2-R_4(z)}{15}
 =\sum_{j=0}^{5}\frac{z^j}{j!}+\frac{z^6}{864}+\frac{z^7}{8640}+\frac{z^8}{138240}.
 \label{eq:richardson_stability}
\end{equation}
Its first exponential defect is
\begin{equation}
 R_{\rm ext}(z)-\e^z=-\frac{z^6}{4320}+O(z^7).
 \label{eq:richardson_defect}
\end{equation}
Thus the affine construction cancels RK4's fifth-order exponential defect and has linear order five. On the one-way zero-dephasing boundary the exact margin vanishes identically, so the same cancellation removes the constituents' leading inward $x^5$ margin and exposes the next defect order. The signed first-defect theorem therefore gives the local physicality failure without a separate Choi expansion:
\begin{equation}
 M_{\rm ext}(x;\varpi)
 =-\frac{G_6(\varpi)}{4320}x^6+O(x^7),
 \qquad
 G_6=\frac{q(q^2-18q+48)}{32},
 \label{eq:richardson_frequency}
\end{equation}
where $q=1+4\varpi^2$. Its two positive orientation thresholds are
\begin{equation}
 \varpi_{\mathrm{ext},-}=\sqrt{2-\frac{\sqrt{33}}4},
 \qquad
 \varpi_{\mathrm{ext},+}=\sqrt{2+\frac{\sqrt{33}}4}.
 \label{eq:richardson_frequency_thresholds}
\end{equation}
The extrapolate is locally outward for $0\le\varpi<\varpi_{\mathrm{ext},-}$, inward between the thresholds, and outward again above $\varpi_{\mathrm{ext},+}$. The Dormand--Prince principal formula has the opposite sequence because it has the same $G_6$ but the opposite sign of $\eta_6$. This all-frequency laboratory-frame orientation remains representation dependent, as does every nonzero-$\varpi$ direct-RK statement. By contrast, the global theorem below is proved at $\varpi=0$, where no rotating-frame change is available or needed; its extrapolation failure is therefore frame independent. Supplemental Fig.~S1 displays the three frequency-orientation bands together.

\begin{theorem}[CPTP loss under Richardson extrapolation of RK4]
\label{thm:richardson}
On nonrotating one-way damping, both $\Phi_{\rm c}(x)$ and $\Phi_{\rm f}(x)$ are CPTP for $0\le x\le\alpha_4$. For every $0<x\le\alpha_4$, the extrapolated map \eqref{eq:richardson_map} is non-CPTP. Its margin factors as
\begin{equation}
 M_{\rm ext}(x)
 =-\frac{x^6P_{10}(x)}{1252412463513600}
 =-\frac{31}{138240}x^6+O(x^7),
 \label{eq:richardson_margin}
\end{equation}
where $P_{10}$ is the degree-ten integer polynomial listed in the Supplementary Material and is strictly positive on $[0,\alpha_4]$. Globally, the extrapolated map later re-enters the CPTP set on the remote interval
\begin{equation}
 [6.034523958649\ldots,6.459127767826\ldots].
 \label{eq:richardson_remote}
\end{equation}
\end{theorem}

\begin{proof}
Equation~\eqref{eq:richardson_defect} and Theorem~\ref{thm:defect} give the leading coefficient in Eq.~\eqref{eq:richardson_margin}. Exact substitution gives the factor $P_{10}$. A Sturm sequence isolates its first positive zero at $6.034523958649\ldots>\alpha_4$, proving outward orientation over the common CPTP interval of the constituents. The complete constraint sign chart then isolates the remote component \eqref{eq:richardson_remote} and its population-ceiling exit.
\end{proof}

The theorem applies when the extrapolated candidate is actually advanced. Many step-doubling codes advance the fine map and use the difference only for error estimation. The general conclusion is candidate-specific: neither the nominal RK method nor the physicality of constituent candidates determines the physicality of a negative-weight affine output.

\begin{figure*}[t]
 \centering
 \makebox[\textwidth][c]{\includegraphics[width=\JCPFigureWidth]{Fig3_candidate_maps.pdf}}
 \caption{Candidate-map complete positivity. (a) At $\varpi=2$, the nontrivial positive-step CPTP components of the full RK4 and two-half-step candidates are disjoint; both sets also contain the zero step. The marked total steps reverse which candidate is CPTP. (b) On nonrotating one-way damping, the coarse and fine candidates are CPTP through $x=\alpha_4$, whereas the negative-weight Richardson extrapolate has a strictly negative Choi margin throughout that common interval.}
 \label{fig:candidate_maps}
\end{figure*}

\section{Executed-candidate certification and adaptive control}
\label{sec:stepcontrol}

Embedded RK and step-doubling methods use local error estimates to propose new steps \cite{DormandPrince1980,HairerNorsettWanner1993,Shampine1985Doubling}. PI and digital-filter controllers refine those proposals using recent error history \cite{Gustafsson1991,Soderlind2002,Soderlind2003}. A complete-positivity failure is not an error measurement. It is a structural rejection of a particular candidate map and should not be inserted into the error history as an artificial error ratio.

\subsection{Executed-input and executed-map semantics}

For a known phase-covariant generator, membership is decided by the four inequalities of Theorem~\ref{thm:criterion} together with the stage-evaluation domain. Reliable implementation is not equivalent to subtracting two nearly equal binary numbers. On the one-way RK4 boundary,
\begin{equation}
 M_4(x;\varpi)=-\frac{x^5q}{147456}\,\mathcal C_q(qx),
\end{equation}
so direct evaluation of $R_4(-x)-|R_4[-x(1/2-\ii\varpi)]|^2$ loses the sign as $x\to0$.

The interface is tri-state:
\begin{equation}
 \textsc{CPTP\_PASS},\qquad \textsc{CPTP\_FAIL},\qquad \textsc{CPTP\_UNCERTAIN}.
 \label{eq:tri_state}
\end{equation}
A Python floating-point input is interpreted as its exact IEEE-754 binary rational, so the certificate concerns the number used by the workflow rather than a nearby printed decimal. Polynomial decisions use exact rational arithmetic and Sturm isolation in the supported cases. An unresolved sign returns \textsc{uncertain}; it is never accepted automatically.

The distinction is observable at the upper RK4 boundary for $\varpi=2$. The printed decimal $x=1.270334626973889$ has margin $+1.43\times10^{-16}$, whereas the exact binary64 value represented by the same display has margin $-9.45\times10^{-17}$. A certifier that rationalizes the printed decimal certifies a different map from the one executed.

Certification is attached to an immutable candidate record, defined here as the binary64 step, construction rule, numerical map, and certificate for the map actually advanced. The estimator, guard, and accumulator pass this record unchanged; local-error estimates remain separate because they compare candidates rather than certify one. After projection, the selected rational step is converted to binary64 and the candidate is recertified at that executed value. Exact-exponential and Strang maps receive a by-construction CPTP certificate only when created internally from a recognized construction rule; direct RK maps are certified algebraically. A caller-supplied label cannot certify an arbitrary matrix.
The certificate applies to the exact mathematical candidate induced by the exact binary input. Its stored complex128 superoperator is generated from the same candidate specification and checked against that specification, but the implementation does not claim a formal enclosure of every downstream BLAS operation. Accordingly, machine-scale negative eigenvalues of the stored matrix near an equality face are monitored realization errors, not substitutes for the exact sign certificate.

\begin{figure*}[t]
 \centering
 \makebox[\textwidth][c]{\includegraphics[width=\JCPFigureWidth]{Fig4_certified_guard.pdf}}
 \caption{Why certification must bind the executed input and candidate. (a) The scaled exact sign polynomial identifies the detached component and the steps skipped by repeated halving. (b) The physical margin scales as $x^5$ and is eventually dominated by fixed absolute tolerances and cancellation in direct binary64 subtraction. Scientific decisions use the exact sign structure or exact candidate constraints, not the displayed floating subtraction.}
 \label{fig:certified_guard}
\end{figure*}

\subsection{One-shot CP projection and error-control separation}

Let $\mathcal H_{\rm CP}^{\rm cand}(h_{\rm prop})$ be the CPTP set of the candidate construction actually proposed over $0<h\le h_{\rm prop}$. For supported phase-covariant polynomial candidates, the certified component locator uses exact root isolation to return the components within this domain. Each root is compared algebraically with both endpoints before its bracket is restricted: exterior roots are discarded, endpoint roots are retained, and interior roots are re-isolated within the domain. An isolated equality point is accepted only if every active inequality and stage-domain condition holds. Any unresolved membership, identity, ordering, bracket overlap, or sign yields \textsc{uncertain}.

\begin{proposition}[Exact feasibility and executed projection]
\label{prop:one_shot_projection}
Assume the component decomposition of $\mathcal H_{\rm CP}^{\rm cand}(h_{\rm prop})$ is certified. Every positive-width component contains an interior rational step specification. The implementation accepts a projection only if it also discovers an IEEE-754 value whose exact binary rational lies in that component and recertifies the corresponding candidate specification. If the finite representable search or binary recertification fails, the status is \textnormal{\textsc{uncertain}} with diagnostic \texttt{EXECUTED\_BINARY\_RECERTIFICATION\_FAILED}. If no positive-width component exists, the routine returns \textnormal{\textsc{no component}}; unresolved root, ordering, membership, or sign information also returns \textnormal{\textsc{uncertain}}. CP feasibility does not determine local accuracy, which must be recomputed for the executed candidate.
\end{proposition}

\begin{proof}
A positive-width component contains a nonempty open interval, hence an interior rational by density of $\mathbb Q$ in $\mathbb R$. This establishes an exact candidate specification but not the success of a finite binary64 search: a narrow component need not contain a value reached by the implementation's representable search stencil. Each proposed binary64 value is therefore interpreted as its exact rational and recertified against every active predicate. Successful recertification proves executed projection; exhaustion or failure leaves the decision unresolved and triggers the stated fail-closed status.
\end{proof}

The controller maintains separate error and CP state. A CP-only rejection leaves the PI error history unchanged while component selection constructs a new candidate. The error estimate is recomputed at that candidate before acceptance. A change of representation or propagator resets the PI memory by policy, because the preceding error ratio no longer belongs to the same candidate family.

\begin{algorithm}[Candidate-aware step]
\label{alg:candidate_controller}
Given a generator, RK method, error-controller proposal $h_{\rm prop}$, candidate-construction rule, tolerances, certification options, and a physical fallback:
\begin{enumerate}
 \item Construct the candidate map intended for advancement and its local-error estimate.
 \item If the error estimate fails, update the ordinary PI/PID proposal and repeat.
 \item Certify that execution record at the same executed numerical input; retain the local-error observation separately.
 \item If the status is \textsc{pass}, accept and update PI memory using that candidate's error.
 \item If the status is \textsc{fail}, isolate the candidate-map CPTP components below the error proposal, choose a certified interior point, reconstruct and recertify the executed candidate there, recompute its error, and return to step 2. If no positive-width component is available, transform or switch propagator, construct the replacement candidate, and return to step 2.
 \item If the status is \textsc{uncertain}, increase precision or construct the physical fallback and return to step 2; never convert uncertainty to acceptance.
\end{enumerate}
\end{algorithm}

At the exact one-way boundary with a detached component, a proposal below the component has no smaller positive direct-RK recovery. With dephasing or bidirectional relaxation, the attached component of Sec.~\ref{sec:robustness} returns, so sufficiently small steps eventually recover; projection versus further reduction is then a cost decision rather than a logical necessity.

\subsection{Adaptive and comparative benchmarks}

The first experiment uses the exact-boundary case $\theta=\kappa=0$, $\varpi=2$, and final dimensionless time $X=2$. Step doubling compares a coarse RK4 map with two half steps and advances the fine map. Four policies are reported over a tolerance sweep: error-only laboratory-frame RK4, a CP guard with rotating-frame fallback, rotating-frame RK4, and Strang composition of exact Hamiltonian and dissipative subflows \cite{Strang1968}. Strang is exact in this commuting benchmark and is included as a structural reference, not as a general timing claim.

Figure~\ref{fig:tolerance_sweep}(a) plots normalized-Choi trace error against a reporting proxy that assigns one unit per RK stage and three per exponential action. This work proxy excludes algebraic certification costs, which are reported separately in Table~\ref{tab:guard_telemetry}. Repeated abscissae indicate equal proxy counts rather than duplicate records: the two error-only observations at work 60 correspond to requested tolerances $10^{-2}$ and $3\times10^{-3}$. At requested tolerance $3\times10^{-4}$, rotating-frame RK4 reaches \BenchmarkRotatingError{} in \BenchmarkRotatingStages{} RK stages; the nearest guard-plus-fallback result reaches \BenchmarkNearestGuardError{} in \BenchmarkNearestGuardStages{} stages.

Across the four boundary rows, every accepted guarded candidate uses rotating-frame fallback. The guard makes \BenchmarkGuardLocatorCalls{} locator calls: \BenchmarkGuardSuccessfulProjections{} projections, \BenchmarkGuardEmptyComponents{} empty component sets, and \BenchmarkGuardUncertainOutcomes{} uncertain results. These runs test fail-closed fallback, not runtime efficiency.

Buffered and near-saddle cases exercise projection at global normalized-Choi errors below $10^{-3}$. At $\theta=0$, $\kappa=10^{-3}$, $\varpi=2$, tolerance $10^{-3}$, and $X=2$, the controller accepts \BenchmarkBufferedDirect{} direct and \BenchmarkBufferedProjected{} projected candidates after \BenchmarkBufferedLocatorCalls{} locator calls, with error \BenchmarkBufferedError{} from \BenchmarkBufferedStages{} RK stages. At $\varpi=1.65$ with the same buffer and tolerance, it makes \BenchmarkNearSaddleProjected{} projection after \BenchmarkNearSaddleLocatorCalls{} locator call, with error \BenchmarkNearSaddleError{}. Figure~\ref{fig:tolerance_sweep}(b) identifies the accepted candidate types; Table~\ref{tab:guard_telemetry} gives the operation counts. A separate large-tolerance path remains in the archive only as an adversarial projection test.

A common dense-qubit noncommuting comparison is also supplied in the archive for the transverse-field generator of Sec.~\ref{sec:robustness}. Direct fixed-step RK4 and exact-subflow Strang are compared against the exact exponential over several subdivisions. This modest baseline illustrates the expected order/structure tradeoff in the same representation; it does not establish a universal ranking against modern low-rank, Kraus, nested-Picard, sparse, or tensor-network implementations.

% Generated mechanically from results/v34/adaptive_benchmark_v34.csv. Do not edit by hand.
\begin{table*}[t]
\centering
\caption{Adaptive safety benchmark. The work proxy counts RK stage evaluations or exponential actions. Algebraic certification costs are excluded and reported separately in Table~\ref{tab:guard_telemetry}.}
\label{tab:adaptive_matrix}
\resizebox{0.99\textwidth}{!}{%
\begin{tabular}{llrrrrrr}
\toprule
Case & Policy & Choi error & Min. eig. & Direct & Projected & Fallback & RK stages\\
\midrule
boundary & error only & $7.92\times 10^{-4}$ & $-1\times 10^{-3}$ & 0 & 0 & 0 & 72\\
boundary & guard + rotating fallback & $7.45\times 10^{-6}$ & $\ge -10^{-15}$ & 0 & 0 & 4 & 132\\
boundary & rotating frame & $1.27\times 10^{-4}$ & $\ge -10^{-15}$ & 0 & 0 & 0 & 36\\
boundary & Strang & $8.33\times 10^{-17}$ & $\ge -10^{-15}$ & 0 & 0 & 0 & 0\\
buffered & guard, buffered & $2.85\times 10^{-4}$ & $\ge -10^{-15}$ & 3 & 2 & 0 & 84\\
buffered & guard, buffered/direct & $2.47\times 10^{-4}$ & $\ge -10^{-15}$ & 5 & 0 & 0 & 60\\
near saddle & guard, near saddle & $2.1\times 10^{-4}$ & $\ge -10^{-15}$ & 1 & 1 & 0 & 48\\
\bottomrule
\end{tabular}%
}
\end{table*}

\begin{table*}[t]
\centering
\caption{Guard telemetry. A locator call is counted on every invocation, including empty and uncertain outcomes. Root-isolation operations count exact Sturm-isolation calls; certification operations count candidate and component certifications, including binary recertification attempts.}
\label{tab:guard_telemetry}
\resizebox{0.99\textwidth}{!}{%
\begin{tabular}{lrrrrrrrrrr}
\toprule
Case & Tol. & Locator & Proj. found & Empty & Uncertain & Root iso. & Cert. & Binary & Prec. esc. & Fallback\\
\midrule
boundary & $1\times 10^{-2}$ & 3 & 0 & 3 & 0 & 9 & 21 & 0 & 0 & 3\\
boundary & $3\times 10^{-3}$ & 3 & 0 & 3 & 0 & 9 & 21 & 0 & 0 & 3\\
boundary & $1\times 10^{-3}$ & 4 & 0 & 4 & 0 & 12 & 33 & 0 & 0 & 4\\
boundary & $3\times 10^{-4}$ & 5 & 0 & 5 & 0 & 15 & 42 & 0 & 0 & 5\\
buffered & $1\times 10^{-3}$ & 2 & 2 & 0 & 0 & 6 & 25 & 2 & 0 & 0\\
buffered & $3\times 10^{-4}$ & 0 & 0 & 0 & 0 & 0 & 15 & 0 & 0 & 0\\
near saddle & $1\times 10^{-3}$ & 1 & 1 & 0 & 0 & 3 & 14 & 1 & 0 & 0\\
\bottomrule
\end{tabular}%
}
\end{table*}


\begin{figure*}[t]
 \centering
 \makebox[\textwidth][c]{\includegraphics[width=\JCPFigureWidth]{Fig5_tolerance_sweep.pdf}}
 \caption{Adaptive safety tests. (a) At $\theta=\kappa=0$, $\varpi=2$, and $X=2$, normalized-Choi trace error is plotted against propagator work. Repeated abscissae are equal-work observations, not duplicate points; the two labeled error-only observations at work 60 use tolerances $10^{-2}$ and $3\times10^{-3}$. The guard curve consists entirely of rotating-frame fallback; Strang is exact for this commuting split and lies at the numerical roundoff floor, so it is omitted from the logarithmic panel. Certification costs are reported separately in Table~\ref{tab:guard_telemetry}. (b) Accepted candidate types in the buffered, near-saddle, and boundary cases. No exact-fallback acceptances occur in these four cases. The boundary guard switches representation because its proposal contains no smaller positive fine-map component.}
 \label{fig:tolerance_sweep}
\end{figure*}

\begin{remark}[Adaptive channel scope]
For each accepted fixed step, the certified candidate is a linear trace-preserving map. If a step sequence is chosen from one particular input state, however, the overall adaptive input--output procedure is generally nonlinear in the initial state and is not itself one reusable quantum channel. Reusable channel construction requires a common generator-, superoperator-, Choi-, or fixed-basis error criterion independent of the state subsequently propagated.
\end{remark}

\section{Operational witness and refinement}
\label{sec:operational}

The Choi test is already a system--ancilla experiment \cite{Choi1975CP,Stinespring1955}. Let $\rho_\Omega=|\Omega\rangle\langle\Omega|/2$ be a maximally entangled two-qubit state. Then
\begin{equation}
 (\id\otimes\Phi)(\rho_\Omega)=\frac{J(\Phi)}{2}.
 \label{eq:bell_choi}
\end{equation}
If $J(\Phi_R)$ has a normalized eigenvector $|v_-\rangle$ with eigenvalue $\lambda_-<0$, the legitimate positive effect $E_-=|v_-\rangle\langle v_-|$ is assigned the numerical Born-rule value
\begin{equation}
 p_-^{\rm RK}=\Tr\!\left[E_-(\id\otimes\Phi_R)(\rho_\Omega)\right]=\frac{\lambda_-}{2}<0.
 \label{eq:negative_born}
\end{equation}
No physical quantum instrument or trajectory ensemble can have such a non-CP map as its unconditional average \cite{DaviesLewis1970,SteinbachGarrawayKnight1995,WonglakhonWisemanChantasri2024}.

For a qubit map, a qubit ancilla already witnesses non-CP. The same witness embeds in every larger ancilla, so the one-qubit admissibility condition is exactly the local safety condition for applying the fixed map as $\Phi_R\otimes\id_d$. This does not claim that direct RK applied to an interacting many-body Liouvillian factorizes into local maps.

Define the normalized-Choi negative mass and Bell-input output trace error by
\begin{equation}
 \delta_{\rm CP}=\frac{\|J_R/2\|_1-1}{2},
 \qquad
 E_J=\frac12\left\|\frac{J_R-J_{\rm ex}}{2}\right\|_1.
 \label{eq:choi_error_definitions}
\end{equation}
Because both normalized Choi operators have trace one, their difference is traceless. The variational characterization of trace distance therefore gives $\delta_{\rm CP}\le E_J$. The ratio can be order one while both quantities are small; it is neither a diamond norm nor a worst-case channel distance.

\subsection{Fixed-horizon convergence from outside the channel set}

Return to nonrotating one-way damping. Fix final dimensionless time $X$ and take $N$ equal direct steps $x_N=X/N$. If
\begin{equation}
 R(z)-\e^z=\eta_mz^m+O(z^{m+1}),
 \qquad
 K=(-1)^m(1-2^{1-m})\eta_m,
\end{equation}
then logarithmic expansion gives
\begin{equation}
 M_N(X)=KXe^{-X}\left(\frac{X}{N}\right)^{m-1}+O(N^{-m}),
 \label{eq:fixed_horizon_margin}
\end{equation}
with the corresponding active Choi eigenvalue
\begin{equation}
 \lambda_{-,N}(X)=\frac{KXe^{-X}}{1+e^{-X}}
 \left(\frac{X}{N}\right)^{m-1}+O(N^{-m}).
 \label{eq:fixed_horizon_eig}
\end{equation}
When $K<0$, the candidate is asymptotically outside the channel set; when $m=p+1$, the negative defect decays at the ordinary global order $p$. A claim of negativity at every finite refinement requires a method-specific sign result, supplied for the displayed Forward Euler and SSPRK(3,3) examples in the Supplementary Material. Thus convergence of maps does not imply physicality of every finite candidate.

\section{Relation to classical diagnostics and structure-preserving alternatives}
\label{sec:comparison}

Complete-positivity stability answers a different question from scalar stability, strong-stability preservation, state-level positivity, or local error control. On one-way damping, the scalar multipliers can both contract while the coupled inequality $|c|^2\le a$ fails. SSP coefficients and absolute-monotonicity radii concern preservation of convex functionals or the positive orthant over classes of ODEs \cite{Kraaijevanger1986Absolute,Kraaijevanger1991Contractivity,GottliebShuTadmor2001,FerracinaSpijker2004,Higueras2005}; positivity-preserving adaptive RK procedures likewise target state or invariant-domain constraints \cite{NussleinRanochaKetcheson2021}. None of these notions is interchangeable with complete positivity of a fixed linear candidate under tensor extension.

The constructive Lindblad literature solves a complementary problem. Kraus and low-rank methods enforce CPTP structure \cite{CaoLu2025,AppeloCheng2025}; nested-Picard schemes address arbitrary order and time-dependent Hamiltonians \cite{HuAppeloCheng2025NPI}; tensor-train and Gregory constructions target compressed or reduced-cost high-order propagation \cite{DelMastroAppeloCheng2026TT,HuAppeloCheng2026Gregory}. The Gregory nested-Picard analysis uses a physics-based test equation to study stability, whereas the finite-step CPTP question here asks whether the particular map executed by an unchanged RK or derived candidate is completely positive. These questions are complementary: test-equation stability does not by itself certify map-level complete positivity. Our contribution is not another universally physical propagator. It describes the exact physicality topology of unchanged RK candidates, including candidates created by composition or extrapolation, and uses that topology to decide whether to accept, transform, project, or switch. Classical extrapolation studies already show that Richardson combinations can change absolute stability, dissipation, and dispersion \cite{Szyszkowicz1985Richardson,ZlatevEtAl2017Richardson,HavasiKazemi2018Richardson}; the result here is the corresponding map-level complete-positivity failure even when both constituent quantum-channel candidates are CPTP.

For a known autonomous phase-covariant qubit generator, one membership test uses two scalar stability-function evaluations and four inequalities, with stage-domain checks when needed. Locating every component or proving emptiness is a separate root-isolation task. The scalar reduction does not extend at the same cost to an arbitrary black-box or higher-dimensional Liouvillian, where a map-level Choi test is substantially larger.

\begin{table*}[t]
\centering
\caption{Roles of the numerical policies used in this paper. The work descriptions are structural and do not assert universal wall-clock rankings across dense, sparse, low-rank, or tensor-compressed representations.}
\label{tab:method_comparison}
\small
\begin{tabularx}{\textwidth}{@{}p{0.15\textwidth}p{0.19\textwidth}p{0.23\textwidth}X@{}}
\toprule
Policy & Candidate advanced & Physicality status & Purpose in this study \\
\midrule
Direct RK4 & unchanged full or fine RK candidate & conditional & baseline with exact generator-conditioned CPTP topology \\
Candidate-aware RK4 & fine RK candidate, projected candidate, or certified fallback & certified conditional & can retain direct RK, select a certified component, or switch; action counts are recorded \\
Rotating-frame RK4 & de-rotated RK candidate followed by exact back-rotation & conditional & representation preconditioning for commuting coherent motion \\
Exact-subflow Strang & symmetric composition of exact CPTP subflows & CPTP by construction & transparent splitting reference; second order when the subflows do not commute \\
Modern Kraus/low-rank/NPI schemes & method-specific structured candidates & enforced in their stated classes & scalable constructive alternatives \cite{CaoLu2025,AppeloCheng2025,HuAppeloCheng2025NPI} \\
\bottomrule
\end{tabularx}
\end{table*}

\subsection{Verification and reproducibility}
\label{subsec:verification}

All analytic identities are regenerated from rational Butcher data. Integer polynomials, Sturm isolation, rational brackets, and joint tests at algebraic roots support the topology and Richardson results; the algebraic procedures follow standard algorithms \cite{BasuPollackRoy2006}. Endpoint membership is decided before bracket restriction, and unresolved domain membership, root identity, ordering, or sign returns \textsc{uncertain}. Direct Choi calculations independently test the scalar criterion on randomized and adversarial candidate families. The companion archive pins the SymPy, NumPy, mpmath, Matplotlib, and pytest environment \cite{MeurerEtAl2017,HarrisEtAl2020,JohanssonMpmath,Hunter2007,Pytest2025}.

All reported values and figures, including the graphical abstract, were regenerated deterministically with the versioned scripts and checked by the procedures described above.

\section{Conclusions}
\label{sec:conclusion}

For the autonomous phase-covariant qubit family, direct RK admits an exact necessary-and-sufficient CPTP criterion. On the laboratory-frame one-way, zero-dephasing ray, classical RK4 exhibits connected, reentrant, isolated, empty-positive-step, and detached regimes that are invisible to scalar stability and ordinary error control.

The principal numerical consequence is candidate dependence. A full RK step, the composition of equal substeps, and a negative-weight extrapolate are different maps over the same total interval. At $|\omega|/\Gamma=2$, the nontrivial positive-step CPTP components of the full-step and two-half-step candidates are disjoint, and the general substep-ladder criterion identifies exactly when consecutive candidate families overlap. Richardson extrapolation raises the linear order of RK4 yet reverses the leading complete-positivity orientation on nonrotating damping; a global certificate shows that the extrapolate remains non-CPTP throughout an interval on which both constituent maps are physical. Adaptive software must therefore certify the map it actually advances.

The guard accepts only a candidate certified at its executed binary64 step. A CP-only rejection leaves the PI error memory unchanged; failed projection, an empty component set, or unresolved certification invokes the physical fallback. Error-only step doubling accepts non-CPTP fine candidates in the boundary test, whereas the guarded workflows do not. Table~\ref{tab:guard_telemetry} reports the component-location and certification workload.

The high-frequency detached topology is representation dependent: an exact rotating frame removes the commuting precession in the central benchmark. In the displayed non-phase-covariant transverse-field family, independent certificates at $\Omega_x/\Gamma\in\{0,0.05,0.10,0.20\}$ each retain two finite-step CP components. These pointwise results do not certify the full interval or arbitrary noncommuting Lindbladians. The one-qubit scalar criterion remains special; time-dependent, higher-dimensional, and interacting generators require map-level certification or structure-preserving propagation.

Complete positivity therefore belongs to the candidate map in its chosen representation, not to the RK method, nominal step, or error estimate. An adaptive code must certify that map or switch to a CPTP construction.


\section*{CRediT authorship contribution statement}
G. Blake Pierpoint: Conceptualization, Methodology, Formal analysis, Investigation, Data curation, Visualization, Project administration, Software, Supervision, Writing -- original draft, Writing -- review and editing. Olivier Bernard: Conceptualization, Methodology, Formal analysis, Investigation, Software, Validation, Writing -- review and editing. Yichen Liu: Software, Methodology, Formal analysis, Investigation, Validation, Writing -- review and editing.

\section*{Acknowledgments}
This research was supported by the Research Computing clusters at Old Dominion University. The Wahab cluster is supported in part by National Science Foundation's grant CNS-1828593, ``MRI Acquisition: A Reconfigurable Computing Infrastructure Enabling Interdisciplinary and Collaborative Research in Hampton Roads.''

\section*{Funding}
This research received no specific grant from funding agencies in the public, commercial, or not-for-profit sectors.

\section*{Declaration of competing interest}
The authors declare that they have no known competing financial interests or personal relationships that could have appeared to influence the work reported in this paper.

\section*{Data availability}
No external data were used. The versioned code-and-data archive contains exact tableaux, symbolic certificates, rational root brackets, interval sign charts, adaptive-benchmark records, plotting data, source code, and regression tests. It will be available from Zenodo at \texttt{[ZENODO DOI TO BE INSERTED BEFORE SUBMISSION]} under the MIT License.

\begin{thebibliography}{99}
\bibitem{Lindblad1976}
G. Lindblad, ``On the generators of quantum dynamical semigroups,'' Commun. Math. Phys. \textbf{48}, 119--130 (1976), \href{https://doi.org/10.1007/BF01608499}{doi:10.1007/BF01608499}.

\bibitem{GoriniKossakowskiSudarshan1976}
V. Gorini, A. Kossakowski, and E. C. G. Sudarshan, ``Completely positive dynamical semigroups of $N$-level systems,'' J. Math. Phys. \textbf{17}, 821--825 (1976), \href{https://doi.org/10.1063/1.522979}{doi:10.1063/1.522979}.

\bibitem{BreuerPetruccione2002}
H.-P. Breuer and F. Petruccione, \textit{The Theory of Open Quantum Systems} (Oxford University Press, Oxford, 2002).

\bibitem{Choi1975CP}
M.-D. Choi, ``Completely positive linear maps on complex matrices,'' Linear Algebra Appl. \textbf{10}, 285--290 (1975), \href{https://doi.org/10.1016/0024-3795(75)90075-0}{doi:10.1016/0024-3795(75)90075-0}.

\bibitem{Kraus1971}
K. Kraus, ``General state changes in quantum theory,'' Ann. Phys. (N.Y.) \textbf{64}, 311--335 (1971), \href{https://doi.org/10.1016/0003-4916(71)90108-4}{doi:10.1016/0003-4916(71)90108-4}.

\bibitem{Stinespring1955}
W. F. Stinespring, ``Positive functions on $C^*$-algebras,'' Proc. Am. Math. Soc. \textbf{6}, 211--216 (1955), \href{https://doi.org/10.1090/S0002-9939-1955-0069403-4}{doi:10.1090/S0002-9939-1955-0069403-4}.

\bibitem{RieschJirauschek2019}
M. Riesch and C. Jirauschek, ``Analyzing the positivity preservation of numerical methods for the Liouville--von Neumann equation,'' J. Comput. Phys. \textbf{390}, 290--296 (2019), \href{https://doi.org/10.1016/j.jcp.2019.04.006}{doi:10.1016/j.jcp.2019.04.006}.

\bibitem{AppeloCheng2025}
D. Appel\"o and Y. Cheng, ``Kraus is king: High-order completely positive and trace preserving low-rank method for the Lindblad master equation,'' J. Comput. Phys. \textbf{534}, 114036 (2025), \href{https://doi.org/10.1016/j.jcp.2025.114036}{doi:10.1016/j.jcp.2025.114036}.

\bibitem{SteinbachGarrawayKnight1995}
J. Steinbach, B. M. Garraway, and P. L. Knight, ``High-order unraveling of master equations for dissipative evolution,'' Phys. Rev. A \textbf{51}, 3302--3308 (1995), \href{https://doi.org/10.1103/PhysRevA.51.3302}{doi:10.1103/PhysRevA.51.3302}.

\bibitem{LeBrisRouchon2013}
C. Le Bris and P. Rouchon, ``Low-rank numerical approximations for high-dimensional Lindblad equations,'' Phys. Rev. A \textbf{87}, 022125 (2013), \href{https://doi.org/10.1103/PhysRevA.87.022125}{doi:10.1103/PhysRevA.87.022125}.

\bibitem{LeBrisRouchonRoussel2015}
C. Le Bris, P. Rouchon, and J. Roussel, ``Adaptive low-rank approximation and denoised Monte Carlo approach for high-dimensional Lindblad equations,'' Phys. Rev. A \textbf{92}, 062126 (2015), \href{https://doi.org/10.1103/PhysRevA.92.062126}{doi:10.1103/PhysRevA.92.062126}.

\bibitem{WernerEtAl2016}
A. H. Werner, D. Jaschke, P. Silvi, M. Kliesch, T. Calarco, J. Eisert, and S. Montangero, ``Positive tensor network approach for simulating open quantum many-body systems,'' Phys. Rev. Lett. \textbf{116}, 237201 (2016), \href{https://doi.org/10.1103/PhysRevLett.116.237201}{doi:10.1103/PhysRevLett.116.237201}.

\bibitem{RieschPiklJirauschek2020}
M. Riesch, A. Pikl, and C. Jirauschek, ``Completely positive trace preserving methods for the Lindblad equation,'' in \textit{2020 International Conference on Numerical Simulation of Optoelectronic Devices (NUSOD)} (IEEE, 2020), pp. 109--110, \href{https://doi.org/10.1109/NUSOD49422.2020.9217670}{doi:10.1109/NUSOD49422.2020.9217670}.

\bibitem{CaoLu2025}
Y. Cao and J. Lu, ``Structure-preserving numerical schemes for Lindblad equations,'' J. Sci. Comput. \textbf{102}, 27 (2025), \href{https://doi.org/10.1007/s10915-024-02707-x}{doi:10.1007/s10915-024-02707-x}.

\bibitem{RobinRouchonSellem2025}
R. Robin, P. Rouchon, and L.-A. Sellem, ``Unconditionally stable time discretization of Lindblad master equations in infinite dimension using quantum channels,'' arXiv:2503.01712 (2025), \href{https://doi.org/10.48550/arXiv.2503.01712}{doi:10.48550/arXiv.2503.01712}.

\bibitem{ChenEtAl2026}
H. Chen, A. Borz\`i, D. Jankovi\'c, J.-G. Hartmann, and P.-A. Hervieux, ``Full- and low-rank exponential Euler integrators for the Lindblad equation,'' SIAM J. Sci. Comput. \textbf{48}, A1--A26 (2026), \href{https://doi.org/10.1137/24M1690795}{doi:10.1137/24M1690795}.

\bibitem{GodinezRamirez2025}
E. Godinez-Ramirez, R. M. Milbradt, and C. B. Mendl, ``Riemannian approach to the Lindbladian dynamics of a locally purified tensor network,'' Phys. Rev. A \textbf{112}, 012224 (2025), \href{https://doi.org/10.1103/kztw-ywxr}{doi:10.1103/kztw-ywxr}.

\bibitem{HuAppeloCheng2025NPI}
J. Hu, D. Appel\"o, and Y. Cheng, ``Arbitrary high order low-rank completely positive and trace preserving (CPTP) schemes for Lindblad equations with time-dependent Hamiltonian,'' arXiv:2511.12012 (2025), \href{https://doi.org/10.48550/arXiv.2511.12012}{doi:10.48550/arXiv.2511.12012}.

\bibitem{DelMastroAppeloCheng2026TT}
P. DelMastro, D. Appel\"o, and Y. Cheng, ``Completely positive and trace preserving schemes with tensor train compression for the Lindblad equation,'' arXiv:2605.01494 (2026), \href{https://doi.org/10.48550/arXiv.2605.01494}{doi:10.48550/arXiv.2605.01494}.

\bibitem{HuAppeloCheng2026Gregory}
J. Hu, D. Appel\"o, and Y. Cheng, ``Gregory nested Picard iteration schemes for open quantum systems governed by the Lindblad equation,'' arXiv:2606.29002 (2026), \href{https://doi.org/10.48550/arXiv.2606.29002}{doi:10.48550/arXiv.2606.29002}.

\bibitem{GuZhang2025}
J. Gu and F. Zhang, ``Generic method for integrating Lindblad master equations,'' Phys. Rev. E \textbf{112}, 044106 (2025), \href{https://doi.org/10.1103/vnhp-fnl3}{doi:10.1103/vnhp-fnl3}.

\bibitem{TureHyeonJang2026}
T. M. Ture, C. Hyeon, and S. J. Jang, ``A simple fourth order propagator based on the Magnus expansion in the Liouville space: Application to a $\Lambda$-system and assessment of the rotating wave approximation,'' J. Chem. Phys. \textbf{164}, 054114 (2026), \href{https://doi.org/10.1063/5.0314773}{doi:10.1063/5.0314773}.

\bibitem{EtienneyRobinRouchon2026}
P.-L. Etienney, R. Robin, and P. Rouchon, ``A posteriori error estimates for the Lindblad master equation,'' Quantum \textbf{10}, 2031 (2026), \href{https://doi.org/10.22331/q-2026-03-16-2031}{doi:10.22331/q-2026-03-16-2031}.

\bibitem{ChenLiLuYing2025}
H. Chen, B. Li, J. Lu, and L. Ying, ``A randomized method for simulating Lindblad equations and thermal state preparation,'' Quantum \textbf{9}, 1917 (2025), \href{https://doi.org/10.22331/q-2025-11-20-1917}{doi:10.22331/q-2025-11-20-1917}.

\bibitem{CaoHeLi2026Unraveling}
Y. Cao, M. He, and X. Li, ``Dynamically optimal unraveling schemes for simulating Lindblad equations,'' J. Phys. A: Math. Theor. \textbf{59}, 165301 (2026), \href{https://doi.org/10.1088/1751-8121/ae5afd}{doi:10.1088/1751-8121/ae5afd}.

\bibitem{BolleyCrouzeix1978}
C. Bolley and M. Crouzeix, ``Conservation de la positivit\'e lors de la discr\'etisation des probl\`emes d'\'evolution paraboliques,'' RAIRO Anal. Num\'er. \textbf{12}, 237--245 (1978), \href{https://doi.org/10.1051/m2an/1978120302371}{doi:10.1051/m2an/1978120302371}.

\bibitem{Kraaijevanger1986Absolute}
J. F. B. M. Kraaijevanger, ``Absolute monotonicity of polynomials occurring in the numerical solution of initial value problems,'' Numer. Math. \textbf{48}, 303--322 (1986), \href{https://doi.org/10.1007/BF01389477}{doi:10.1007/BF01389477}.

\bibitem{Kraaijevanger1991Contractivity}
J. F. B. M. Kraaijevanger, ``Contractivity of Runge--Kutta methods,'' BIT Numer. Math. \textbf{31}, 482--528 (1991), \href{https://doi.org/10.1007/BF01933264}{doi:10.1007/BF01933264}.

\bibitem{Horvath1998Positivity}
Z. Horv\'ath, ``Positivity of Runge--Kutta and diagonally split Runge--Kutta methods,'' Appl. Numer. Math. \textbf{28}, 309--326 (1998), \href{https://doi.org/10.1016/S0168-9274(98)00050-6}{doi:10.1016/S0168-9274(98)00050-6}.

\bibitem{Horvath2005Threshold}
Z. Horv\'ath, ``On the positivity step size threshold of Runge--Kutta methods,'' Appl. Numer. Math. \textbf{53}, 341--356 (2005), \href{https://doi.org/10.1016/j.apnum.2004.08.026}{doi:10.1016/j.apnum.2004.08.026}.

\bibitem{GottliebShuTadmor2001}
S. Gottlieb, C.-W. Shu, and E. Tadmor, ``Strong stability-preserving high-order time discretization methods,'' SIAM Rev. \textbf{43}, 89--112 (2001), \href{https://doi.org/10.1137/S003614450036757X}{doi:10.1137/S003614450036757X}.

\bibitem{FerracinaSpijker2004}
L. Ferracina and M. N. Spijker, ``Stepsize restrictions for the total-variation-diminishing property in general Runge--Kutta methods,'' SIAM J. Numer. Anal. \textbf{42}, 1073--1093 (2004), \href{https://doi.org/10.1137/S0036142902415584}{doi:10.1137/S0036142902415584}.

\bibitem{Higueras2005}
I. Higueras, ``Representations of Runge--Kutta methods and strong stability preserving methods,'' SIAM J. Numer. Anal. \textbf{43}, 924--948 (2005), \href{https://doi.org/10.1137/S0036142903427068}{doi:10.1137/S0036142903427068}.

\bibitem{BlanesIserlesMacnamara2022}
S. Blanes, A. Iserles, and S. Macnamara, ``Positivity-preserving methods for ordinary differential equations,'' ESAIM Math. Model. Numer. Anal. \textbf{56}, 1843--1870 (2022), \href{https://doi.org/10.1051/m2an/2022042}{doi:10.1051/m2an/2022042}.

\bibitem{FujiwaraAlgoet1999}
A. Fujiwara and P. Algoet, ``One-to-one parametrization of quantum channels,'' Phys. Rev. A \textbf{59}, 3290--3294 (1999), \href{https://doi.org/10.1103/PhysRevA.59.3290}{doi:10.1103/PhysRevA.59.3290}.

\bibitem{RuskaiSzarekWerner2002}
M. B. Ruskai, S. Szarek, and E. Werner, ``An analysis of completely-positive trace-preserving maps on $M_2$,'' Linear Algebra Appl. \textbf{347}, 159--187 (2002), \href{https://doi.org/10.1016/S0024-3795(01)00547-X}{doi:10.1016/S0024-3795(01)00547-X}.

\bibitem{BraunEtAl2014}
D. Braun, O. Giraud, I. Nechita, C. Pellegrini, and M. \v{Z}nidari\v{c}, ``A universal set of qubit quantum channels,'' J. Phys. A: Math. Theor. \textbf{47}, 135302 (2014), \href{https://doi.org/10.1088/1751-8113/47/13/135302}{doi:10.1088/1751-8113/47/13/135302}.

\bibitem{Hall2008Complete}
M. J. W. Hall, ``Complete positivity for time-dependent qubit master equations,'' J. Phys. A: Math. Theor. \textbf{41}, 205302 (2008), \href{https://doi.org/10.1088/1751-8113/41/20/205302}{doi:10.1088/1751-8113/41/20/205302}.

\bibitem{Lankinen2016Complete}
J. Lankinen, H. Lyyra, B. Sokolov, J. Teittinen, B. Ziaei, and S. Maniscalco, ``Complete positivity, finite-temperature effects, and additivity of noise for time-local qubit dynamics,'' Phys. Rev. A \textbf{93}, 052103 (2016), \href{https://doi.org/10.1103/PhysRevA.93.052103}{doi:10.1103/PhysRevA.93.052103}.

\bibitem{Filippov2020Phase}
S. N. Filippov, A. N. Glinov, and L. Lepp\"aj\"arvi, ``Phase covariant qubit dynamics and divisibility,'' Lobachevskii J. Math. \textbf{41}, 617--630 (2020), \href{https://doi.org/10.1134/S1995080220040095}{doi:10.1134/S1995080220040095}.

\bibitem{Siudzinska2022Mixtures}
K. Siudzi\'nska, ``Phase-covariant mixtures of non-unital qubit maps,'' J. Phys. A: Math. Theor. \textbf{55}, 405303 (2022), \href{https://doi.org/10.1088/1751-8121/ac909b}{doi:10.1088/1751-8121/ac909b}.

\bibitem{Siudzinska2023Geometry}
K. Siudzi\'nska, ``Geometry of phase-covariant qubit channels,'' J. Phys. Commun. \textbf{7}, 075002 (2023), \href{https://doi.org/10.1088/2399-6528/ace0f4}{doi:10.1088/2399-6528/ace0f4}.

\bibitem{LonigroChruscinski2023}
D. Lonigro and D. Chru\'sci\'nski, ``Excitation-damping quantum channels,'' J. Phys. A: Math. Theor. \textbf{56}, 255301 (2023), \href{https://doi.org/10.1088/1751-8121/acd734}{doi:10.1088/1751-8121/acd734}.

\bibitem{Shampine1985Doubling}
L. F. Shampine, ``Local error estimation by doubling,'' Computing \textbf{34}, 179--190 (1985), \href{https://doi.org/10.1007/BF02259844}{doi:10.1007/BF02259844}.

\bibitem{NussleinRanochaKetcheson2021}
S. N\"u{\ss}lein, H. Ranocha, and D. I. Ketcheson, ``Positivity-preserving adaptive Runge--Kutta methods,'' Commun. Appl. Math. Comput. Sci. \textbf{16}, 155--179 (2021), \href{https://doi.org/10.2140/camcos.2021.16.155}{doi:10.2140/camcos.2021.16.155}.

\bibitem{Szyszkowicz1985Richardson}
M. Szyszkowicz, ``The intervals of stability of Runge--Kutta methods after Richardson extrapolation,'' Appl. Math. (Warsaw) \textbf{18}, 657--661 (1985), \href{https://doi.org/10.4064/am-18-4-657-661}{doi:10.4064/am-18-4-657-661}.

\bibitem{ZlatevEtAl2017Richardson}
Z. Zlatev, I. Dimov, I. Farag\'o, K. Georgiev, and \`A. Havasi, ``Stability of the Richardson extrapolation combined with some implicit Runge--Kutta methods,'' J. Comput. Appl. Math. \textbf{310}, 224--240 (2017), \href{https://doi.org/10.1016/j.cam.2016.03.018}{doi:10.1016/j.cam.2016.03.018}.

\bibitem{HavasiKazemi2018Richardson}
\`A. Havasi and E. Kazemi, ``On Richardson extrapolation for low-dissipation low-dispersion diagonally implicit Runge--Kutta schemes,'' J. Comput. Phys. \textbf{358}, 21--35 (2018), \href{https://doi.org/10.1016/j.jcp.2017.12.043}{doi:10.1016/j.jcp.2017.12.043}.

\bibitem{PuzzuoliEtAl2023}
D. Puzzuoli, C. J. Wood, D. J. Egger, B. Rosand, and K. Ueda, ``Qiskit Dynamics: A Python package for simulating the time dynamics of quantum systems,'' J. Open Source Softw. \textbf{8}, 5853 (2023), \href{https://doi.org/10.21105/joss.05853}{doi:10.21105/joss.05853}.

\bibitem{HairerNorsettWanner1993}
E. Hairer, S. P. N\o rsett, and G. Wanner, \textit{Solving Ordinary Differential Equations I: Nonstiff Problems}, 2nd ed. (Springer, Berlin, 1993), \href{https://doi.org/10.1007/978-3-540-78862-1}{doi:10.1007/978-3-540-78862-1}.

\bibitem{IserlesNorsett1991OrderStars}
A. Iserles and S. P. N\o rsett, \textit{Order Stars: Theory and Applications} (Chapman \& Hall, London, 1991).

\bibitem{DormandPrince1980}
J. R. Dormand and P. J. Prince, ``A family of embedded Runge--Kutta formulae,'' J. Comput. Appl. Math. \textbf{6}, 19--26 (1980), \href{https://doi.org/10.1016/0771-050X(80)90013-3}{doi:10.1016/0771-050X(80)90013-3}.

\bibitem{DaviesLewis1970}
E. B. Davies and J. T. Lewis, ``An operational approach to quantum probability,'' Commun. Math. Phys. \textbf{17}, 239--260 (1970), \href{https://doi.org/10.1007/BF01647093}{doi:10.1007/BF01647093}.

\bibitem{WonglakhonWisemanChantasri2024}
N. Wonglakhon, H. M. Wiseman, and A. Chantasri, ``Completely positive trace-preserving maps for higher-order unraveling of Lindblad master equations,'' Phys. Rev. A \textbf{110}, 062207 (2024), \href{https://doi.org/10.1103/PhysRevA.110.062207}{doi:10.1103/PhysRevA.110.062207}.

\bibitem{Gustafsson1991}
K. Gustafsson, ``Control theoretic techniques for stepsize selection in explicit Runge--Kutta methods,'' ACM Trans. Math. Softw. \textbf{17}, 533--554 (1991), \href{https://doi.org/10.1145/210232.210242}{doi:10.1145/210232.210242}.

\bibitem{Soderlind2002}
G. S\"oderlind, ``Automatic control and adaptive time-stepping,'' Numer. Algorithms \textbf{31}, 281--310 (2002), \href{https://doi.org/10.1023/A:1021160023092}{doi:10.1023/A:1021160023092}.

\bibitem{Soderlind2003}
G. S\"oderlind, ``Digital filters in adaptive time-stepping,'' ACM Trans. Math. Softw. \textbf{29}, 1--26 (2003), \href{https://doi.org/10.1145/641876.641877}{doi:10.1145/641876.641877}.

\bibitem{Strang1968}
G. Strang, ``On the construction and comparison of difference schemes,'' SIAM J. Numer. Anal. \textbf{5}, 506--517 (1968), \href{https://doi.org/10.1137/0705041}{doi:10.1137/0705041}.

\bibitem{BasuPollackRoy2006}
S. Basu, R. Pollack, and M.-F. Roy, \textit{Algorithms in Real Algebraic Geometry}, 2nd ed. (Springer, Berlin, 2006), \href{https://doi.org/10.1007/3-540-33099-2}{doi:10.1007/3-540-33099-2}.

\bibitem{MeurerEtAl2017}
A. Meurer, C. P. Smith, M. Paprocki, O. \v{C}ert\'ik, S. B. Kirpichev, M. Rocklin, A. Kumar, S. Ivanov, J. K. Moore, S. Singh, T. Rathnayake, S. Vig, B. E. Granger, R. P. Muller, F. Bonazzi, H. Gupta, S. Vats, F. Johansson, F. Pedregosa, M. J. Curry, A. R. Terrel, \v{S}. Rou\v{c}ka, A. Saboo, I. Fernando, S. Kulal, R. Cimrman, and A. Scopatz, ``SymPy: symbolic computing in Python,'' PeerJ Comput. Sci. \textbf{3}, e103 (2017), \href{https://doi.org/10.7717/peerj-cs.103}{doi:10.7717/peerj-cs.103}.

\bibitem{HarrisEtAl2020}
C. R. Harris, K. J. Millman, S. J. van der Walt, R. Gommers, P. Virtanen, D. Cournapeau, E. Wieser, J. Taylor, S. Berg, N. J. Smith, R. Kern, M. Picus, S. Hoyer, M. H. van Kerkwijk, M. Brett, A. Haldane, J. F. del R\'io, M. Wiebe, P. Peterson, P. G\'erard-Marchant, K. Sheppard, T. Reddy, W. Weckesser, H. Abbasi, C. Gohlke, and T. E. Oliphant, ``Array programming with NumPy,'' Nature \textbf{585}, 357--362 (2020), \href{https://doi.org/10.1038/s41586-020-2649-2}{doi:10.1038/s41586-020-2649-2}.

\bibitem{JohanssonMpmath}
F. Johansson and the mpmath development team, \textit{mpmath: a Python library for arbitrary-precision floating-point arithmetic}, version 1.3.0 (2023), \href{https://mpmath.org/}{https://mpmath.org/}.

\bibitem{Hunter2007}
J. D. Hunter, ``Matplotlib: A 2D graphics environment,'' Comput. Sci. Eng. \textbf{9}, 90--95 (2007), \href{https://doi.org/10.1109/MCSE.2007.55}{doi:10.1109/MCSE.2007.55}.

\bibitem{Pytest2025}
pytest development team, \textit{pytest 9.0.2 documentation} (2025), \href{https://docs.pytest.org/en/9.0.x/}{https://docs.pytest.org/en/9.0.x/}.

\end{thebibliography}

